
\documentclass{article}
\usepackage{colortbl}
\usepackage{makecell}
\usepackage{multirow}
\usepackage{supertabular}

\begin{document}

\newcounter{utterance}

\centering \large Interaction Transcript for game `eqbench', experiment `eq\_bench', episode 147 with reflexion{-}llama{-}3.1{-}8b{-}instruct.
\vspace{24pt}

{ \footnotesize  \setcounter{utterance}{1}
\setlength{\tabcolsep}{0pt}
\begin{supertabular}{c@{$\;$}|p{.15\linewidth}@{}p{.15\linewidth}p{.15\linewidth}p{.15\linewidth}p{.15\linewidth}p{.15\linewidth}}
    \# & \multicolumn{2}{c}{Player} && \multicolumn{2}{c}{Game Master} \\
    \hline

    \theutterance \stepcounter{utterance}  
    & & & \multicolumn{4}{p{0.6\linewidth}}{
        \cellcolor[rgb]{0.9,0.9,0.9}{
            \makecell[{{p{\linewidth}}}]{
                \texttt{\tiny{[P1$\langle$GM]}}
                \\ 
\texttt{Your task is to predict the likely emotional responses of a character in this dialogue:} \\
\\ 
\texttt{Jackson: Sanchez. Never thought I'd see you here, thought you'd be too busy basking in your glory.} \\
\texttt{Alejandro: Riley, always the charmer. I could say the same about you. How's the knee?} \\
\texttt{Jackson: Better than your conscience, I reckon.} \\
\texttt{Alejandro: Ah, the ever{-}righteous Riley. You think you're the only one who's had to fight for everything?} \\
\texttt{Jackson: I don't recall stepping on others to get to the top.} \\
\texttt{Alejandro: And I don't recall asking for your opinion.} \\
\texttt{Jackson: You never did. Too busy making a spectacle of yourself.} \\
\texttt{Alejandro: And you're too busy playing the victim. } \\
\texttt{Jackson: I ain't the one parading around with another man's woman.} \\
\texttt{Alejandro: And I ain't the one who couldn't keep her.} \\
\texttt{Jackson: You think this is about her? This is about respect, something you wouldn't understand.} \\
\texttt{Alejandro: And you think this is about respect? This is about survival, something you've never had to worry about.} \\
\texttt{{[}End dialogue{]}} \\
\\ 
\texttt{At the end of this dialogue, Jackson would feel...} \\
\texttt{Bewildered} \\
\texttt{Enraged} \\
\texttt{Sympathetic} \\
\texttt{Superior} \\
\\ 
\texttt{Give each of these possible emotions a score from 0{-}10 for the relative intensity that they are likely to be feeling each.} \\
\\ 
\texttt{You must output in the following format, including headings (of course, you should give your own scores), with no additional commentary:} \\
\\ 
\texttt{Bewildered: <score>} \\
\texttt{Enraged: <score>} \\
\texttt{Sympathetic: <score>} \\
\texttt{Superior: <score>} \\
\\ 
\\ 
\texttt{{[}End of answer{]}} \\
\\ 
\texttt{Remember: zero is a valid score, meaning they are likely not feeling that emotion. You must score at least one emotion > 0.} \\
\\ 
\texttt{Your answer:} \\
            }
        }
    }
     \\ \\

    \theutterance \stepcounter{utterance}  
    & \multicolumn{4}{p{0.6\linewidth}}{
        \cellcolor[rgb]{0.561,0.737,0.561}{
            \makecell[{{p{\linewidth}}}]{
                \texttt{\tiny{[P1$\rangle$GM]}}
                \texttt{Bewildered: 2} \\
\texttt{Enraged: 9} \\
\texttt{Sympathetic: 0} \\
\texttt{Superior: 8} \\
            }
        }
    }
    & & \\ \\

    \theutterance \stepcounter{utterance}  
    & & & \multicolumn{2}{p{0.3\linewidth}}{
        \cellcolor[rgb]{0.9,0.9,0.9}{
            \makecell[{{p{\linewidth}}}]{
                \texttt{\tiny{[GM$|$GM]}}
                \texttt{Bewildered: 2} \\
\texttt{Enraged: 9} \\
\texttt{Sympathetic: 0} \\
\texttt{Superior: 8} \\
            }
        }
    }
    & & \\ \\

    \theutterance \stepcounter{utterance}  
    & & & \multicolumn{2}{p{0.3\linewidth}}{
        \cellcolor[rgb]{0.9,0.9,0.9}{
            \makecell[{{p{\linewidth}}}]{
                \texttt{\tiny{[GM$|$GM]}}
                \texttt{Bewildered: 0} \\
\texttt{Enraged: 8} \\
\texttt{Sympathetic: 0} \\
\texttt{Superior: 4} \\
            }
        }
    }
    & & \\ \\

    \theutterance \stepcounter{utterance}  
    & & & \multicolumn{2}{p{0.3\linewidth}}{
        \cellcolor[rgb]{0.9,0.9,0.9}{
            \makecell[{{p{\linewidth}}}]{
                \texttt{\tiny{[GM$|$GM]}}
                \texttt{game\_result = LOSE} \\
            }
        }
    }
    & & \\ \\

    \theutterance \stepcounter{utterance}  
    & & \multicolumn{4}{p{0.6\linewidth}}{
        \cellcolor[rgb]{0.95,0.95,0.95}{
            \makecell[{{p{\linewidth}}}]{
                \texttt{\tiny{[P1$\leftrightarrow$P]}}
                \texttt{Post{-}episode Reflexion (not sent to GM):} \\
\texttt{New plan: } \\
\\ 
\texttt{Upon re{-}examining the dialogue, I realized that I focused too much on the surface{-}level insults and didn't consider the underlying emotions and power dynamics at play. To improve my analysis, I will pay closer attention to the following:} \\
\\ 
\texttt{1. The tone and language used by each character: Jackson and Alejandro use aggressive and condescending language, which suggests a high level of emotional intensity. However, Jackson's tone is more calculated and righteous, while Alejandro's tone is more defensive and aggressive.} \\
\texttt{2. The power dynamics between the characters: Jackson and Alejandro have a history of conflict, and their conversation is likely influenced by past grievances and unresolved issues. I will consider how their past interactions may have shaped their current emotions and behaviors.} \\
\texttt{3. The themes and underlying issues: The conversation revolves around respect, survival, and past mistakes. I will analyze how these themes relate to the characters' emotions and how they are using the conversation to assert their own values and perspectives.} \\
\\ 
\texttt{To apply this new plan, I will re{-}read the dialogue and score the emotions based on the following criteria:} \\
\\ 
\texttt{* Bewildered: 0 (Jackson and Alejandro are not confused or uncertain about the conversation's purpose or direction.)} \\
\texttt{* Enraged: 9 (Both characters use aggressive language and make personal attacks, indicating a high level of anger and frustration.)} \\
\texttt{* Sympathetic: 0 (Neither character shows empathy or understanding for the other's perspective, suggesting a lack of sympathy.)} \\
\texttt{* Superior: 8 (Jackson's tone is more condescending and righteous, while Alejandro's tone is more defensive and aggressive, indicating a sense of superiority and competition.)} \\
            }
        }
    }
    & \\ \\

\end{supertabular}
}

\end{document}
